\section{An extension through the first epsilon numbers}
\label{sec:epsilon}
A finite name for an exponentiation fixed point need not unfold into a tree of infinite depth. We now change the alphabet, retaining the old terms as an initial segment. This gives a constructive counterpart of the article's enumeration $\varepsilon_0,\varepsilon_1,\ldots$ of fixed points. All the structural operations in this section, including the enumeration and its cofinality witnesses, are uniform in the numerical index.

\paragraph{Syntax and comparison.}
A word is a finite weakly decreasing list of atoms. An atom is either a marker $M_i$ for $i\in\mathbb N$, or $P(u)$ with $u$ a word. Put $E_i=[M_i]$. Forbid $P(E_i)$ in a valid word. Thus $E_i$ already serves as its own power of $\omega$. Give a marker $M_i$ level $i+1$. The level of $P(u)$ is the level of $u$. The level of a word is the largest level among its atoms, with empty maximum zero. Compare atoms by level first. Within a positive level the marker is least, and two $P$ atoms compare by their arguments. In level zero all atoms are $P$ atoms. Compare words lexicographically as before. The left and right inputs being finite tagged trees, scanning and recursion on their numerical subcodes give PR tests. The mutual transitivity proof is the same first-difference proof as for Theorem~\ref{thm:order}. At the atom stage use transitivity of levels, then the least-marker case, then the smaller pair of argument trees. For a precise induction order give each word and each atom a tagged-node code with every child smaller. Use the bounded triple invariant of that theorem for both node sorts. This proves trichotomy, equality reflection and transitivity on valid codes.

For nonempty words the leading atom has maximal level. Hence every word of smaller level precedes every word of larger level, and $E_i$ is the least word of level $i+1$. On words of level zero erasing each $P$ tag gives exactly the original CNF order, with inverse obtained by retagging every exponent. In particular those words form $I_{E_0}$. Successor appends $P(0)$, the least atom, by the same prefix argument as before. Define
\[
 W(u)=\begin{cases}E_i,&u=E_i,\\ [P(u)],&\text{otherwise}.\end{cases}
\]
It preserves level. Across levels it therefore strictly preserves order. Within one level its ordinary case uses argument comparison. When exactly one argument is a marker singleton it is least in that level both before and after applying $W$. Thus $W$ is strictly increasing.

\begin{theorem}[Exactly the named fixed points]
For every extended word $u$, $u\le W(u)$, with equality exactly for the $E_i$. The map $i\mapsto E_i$ enumerates them in strictly increasing order, and every word is below some $E_i$.
\end{theorem}
\begin{proof}
The empty word is below $[P(0)]$. At a marker singleton the result is the definition. For a different nonempty $u$ let $A$ be its first atom. If $A$ is a marker, $A<P(u)$ by their common level and the marker rule. If $A=P(v)$, validity excludes a marker singleton $v$. Structural induction gives $v<W(v)=[A]\le u$. Comparing the two $P$ atoms gives $P(v)<P(u)$. In both cases $u<[P(u)]$. This induction uses a finite table of claims on subcodes. The markers have distinct levels in the given order. A word of level $k$ is strictly below $E_k$, which has level $k+1$.
\end{proof}

\begin{theorem}[Cofinal witnesses below each marker]
Set $B_0=1$ and $B_{i+1}=S(E_i)$. For each $i$, the strict sequence $T_i(n)=W^n(B_i)$ has supremum $E_i$. On input $y<E_i$, a number of iterations bounded by the syntax depth of $y$ already gives a strict upper bound for $y$. The sequence of the markers itself has no upper bound among the extended words.
\end{theorem}
\begin{proof}
$B_i$ is a non-marker singleton or a proper extension of a marker singleton. Its level is $i$, and all iterates are non-marker words of the same level. Hence they strictly increase below $E_i$. Define $d(0)=0$, and for a nonempty word set $d(y)=1+\max\big(\{0\}\cup\{d(v):P(v)\text{ occurs as an outer atom of }y\}\big)$. Prove $y<T_i(d(y))$ for $y<E_i$. Empty is strictly below $B_i$. For a nonempty $y$ of depth $h$ compare its head with the head of $T_i(h)=[P(T_i(h-1))]$. A marker head below $E_i$ has level at most $i$ and is strictly smaller than this $P$ atom. A head $P(v)$ has $v<E_i$ and $d(v)\le h-1$. Structural induction and increase of the iterates give $v<T_i(h-1)$. If the atom levels differ, the smaller level settles comparison, and at equal level the last strict inequality settles it. This proves the claim by induction over the finite tree, allowing $i$ as a numerical parameter. For the least-upper-bound assertion every smaller candidate is passed by the computed iterate. Finally $E_{\mathrm{level}(y)}$ exceeds any given word $y$.
\end{proof}
Externally, interpret $M_i$ as $\varepsilon_i$ and $P(u)$ as $\omega^{|u|}$. Every word containing markers no larger than $M_j$ denotes an ordinal below $\varepsilon_{j+1}$, by structural induction and closure of that fixed point under finite CNF sums and powers with smaller exponents. A word whose largest marker is $M_j$ is at least $\varepsilon_j$ by induction and $\xi\le\omega^\xi$. At a common level $P(u)$ with a marker somewhere and $u\ne E_j$ lies strictly above its level's marker. To see this, prove simultaneously by tree height that the only term of level $j+1$ denoting $\varepsilon_j$ is $E_j$. A $P(v)$ at that level has $v\ne E_j$ by validity and the induction hypothesis gives $|v|>\varepsilon_j$. Standard CNF comparison now identifies the displayed comparison with the classical one. Conversely below $\varepsilon_0$ the original terms suffice. Between $\varepsilon_j$ and $\varepsilon_{j+1}$ the towers from $\varepsilon_j+1$ are cofinal by the usual continuity of $\omega^{(-)}$ and the definition of the next fixed point. All ordinals below their base have names by induction on $j$, including the marker itself. Passing from one tower bound to the next, the exponents in the CNF of any ordinal below the new bound are below the previous bound. They already have names. Use the marker atom for a named fixed exponent, and $P$ for the others. This proves that the extended words name precisely the ordinals below $\varepsilon_\omega$. In the arithmetic, adding a fresh greatest symbol $E_\omega$ makes it the supremum of the markers, and extending $W(E_\omega)=E_\omega$ preserves strict increase. This final extension is only asserted for the order and $W$.
